% !TeX root = ../papers/nc-p3.tex
\section{量子部分の全特殊値を含むindexed同型判定}
\label{sec:quantum-indexed-all}

$q_{ba}=q_{ab}^{-1}$, $q_{aa}=1$と拡張し、異なる三添字に対して
$T_{abc}(q)=q_{ab}q_{bc}q_{ca}$と置く。
\begin{theorem}[量子多項式の全パラメータindexed同型]\label{thm:quantum-indexed-all}
全てのパラメータが非零である量子多項式代数$A(q),A(q')$について、次は同値である。
\begin{enumerate}
\item $\check A(q)$と$\check A(q')$が添字保持同型である。
\item ある置換$\sigma\in S_4$と$d_0,\ldots,d_3\in k^\times$が存在し、
全ての異なる$a,b$について
$q'_{\sigma(a)\sigma(b)}=q_{ab}d_a/d_b$となる。
\item ある$\sigma\in S_4$について、全ての三添字で
$T_{abc}(q)=T_{\sigma(a)\sigma(b)\sigma(c)}(q')$となる。
\end{enumerate}
この判定は根が1となる特殊値やパラメータの重複も含み、一般性条件を要しない。
これは量子多項式部分を尽くす同型条件であり、全AS四窓の分類ではない。
\end{theorem}
\begin{proof}
(2)なら位置$i$で$g_i(x_a)=d_a^{-i}x'_{\sigma(a)}$と置く。
二次関係の両項の係数比を取れば、全整数$i$で関係空間が対応し、(1)を得る。
(2)から(3)は三つの比の相殺による。
逆に(3)があれば、置換後の比$r_{ab}=q'_{\sigma(a)\sigma(b)}/q_{ab}$に対し
$r_{ab}r_{bc}=r_{ac}$である。$d_0=1$, $d_b=r_{0b}^{-1}$と取れば$r_{ab}=d_a/d_b$となる。

残る(1)から(2)を示す。
最初の二次関係から得る点対応は
\[
 x_by_a=q_{ab}x_ay_b\qquad(a<b)
\]
である。$x$の非零座標集合を$I$とすると、非零$y$が存在するための必要十分条件は
$I$内の全ての三添字について$T_{abc}=1$となることである。
実際、$y$の非零座標も$I$と一致し、$y_b/x_b=q_{ab}^{-1}(y_a/x_a)$の整合性がこの条件になる。
その場合の$y$は射影的に一意であり、特に各座標点$p_a$の対応先は$p_a$自身である。
従って左射影支持は、全三添字条件を満たす$I$の座標線型空間$\mathbb P(I)$の和となる。

四つの三添字について
$T_{012}T_{023}=T_{013}T_{123}$なので、$T=1$となる三角形の数は$0,1,2,4$のいずれかである。
左支持は順に、六本の座標直線、一平面と外の一点から出る三本の直線、
二平面とそれぞれの外の頂点を結ぶ一本の直線、全$\mathbb P^3$となる。
添字保持同型の最初の二つの行列を$g_0,g_1$とすると、これらは点対応の両因子の射影変換を与える。

三角形数が0の場合、六直線の三重点として四座標点が内在的に復元される。
数が1の場合も、唯一の平面上の三つの付着点と、外の三直線の共通交点として復元される。
従ってこれらの場合$g_0$は単項行列である。
各座標点の右対応が一意なので$g_1$も同じ置換を持つ単項行列となる。
対角係数を$a_j,b_j$と書くと、関係保存から
$q'_{\sigma(a)\sigma(b)}=q_{ab}(a_a/b_a)/(a_b/b_b)$を得る。

三角形数が2の場合は、添字を付け直して共通の辺を$01$、外の二頂点を$2,3$とする。
二平面の交線と、残る直線の二つの付着点が内在的であるから、
$g_0$は$\langle x_0,x_1\rangle$上の$\GL_2$ブロックと外の二頂点のスカラーからなる。
まず(2)から(1)という既に示した基底変換を使い、双方のパラメータで$q_{01}=1$に正規化してよい。
$T_{012}=T_{013}=1$から$q_{02}=q_{12}$, $q_{03}=q_{13}$となる。
従って$\langle x_0,x_1\rangle$内の任意の$\GL_2$変換は量子代数のgraded自己同型である。
これを相手側で合成して$g_0$の$\GL_2$ブロックを除けば、前と同じ単項行列の場合に帰着する。
この正規化と自己同型の除去は、結論(2)の成立を変えない。

三角形数が4なら、双方の全ての$T$が1であるから、(3)から(2)という既証部分を直接適用できる。
以上で全特殊値を尽くした。
\end{proof}
\begin{remark}
三角形数0の一般量子対象で得られる従来の3パラメータの忘却を、
上の証明は平面成分とその交線に$\GL_2$自由度が現れる特殊値まで拡張する。
$\rho$の一致はこの三添字判定より弱い。
同じ$\rho=8$でも、$q_{01}=q_{12}=q_{23}=1$で
$(q_{02},q_{13},q_{03})=(2,2,1)$と$(3,3,1)$の二対象は非同型である。
\path{computations/verify_quantum_indexed.py}は24置換の判定と対角証明書を検査する。
\end{remark}
\begin{corollary}[量子部分の三変数正規形と有限残余作用]\label{cor:quantum-torus-normal-form}
量子多項式の全indexed同型類は、
$q_{01}=q_{02}=q_{03}=1$, $(q_{12},q_{13},q_{23})=(p,q,r)\in(k^\times)^3$
という三変数表示で被覆される。
二つの表示が同型であることは、次の三つの変換が生成する$S_4$作用で同じ軌道に属することと同値である。
\[
 s_{01}(p,q,r)=(p^{-1},q^{-1},pr/q),\qquad
 s_{12}(p,q,r)=(p^{-1},r,q),\qquad
 s_{23}(p,q,r)=(q,p,r^{-1}).
\]
従って量子部分の残余同型は高々24個の明示的候補で判定でき、特殊値もこの式で扱える。
これは粗い同型類の記述であり、自己同型群の次元が特殊値でも一定という主張ではない。
\end{corollary}
\begin{proof}
$d_0=1$, $d_b=q_{0b}$という前定理の変換で最初の三係数を1へ直せる。
残る三係数は$(T_{012},T_{013},T_{023})$であり、第四の三添字値は$T_{123}=pr/q$である。
置換$(01),(12),(23)$の作用をこの三添字値に適用すれば上の三式になる。
前定理が残余同型の必要十分性を与える。
\end{proof}
